\section{灵敏度分析}

选取突水流量、第二水源启动时刻、矿工速度和安全水深阈值进行单因素扰动。设参数 $p$ 的扰动为
$\Delta p$，模型输出变化为 $\Delta F$，相对灵敏度定义为
\begin{equation}
  S_p = \frac{\Delta F/F}{\Delta p/p}.
\end{equation}

流量在 $\pm10\%$ 范围内变化时，两矿井覆盖数量不变，最大到达和充满时刻与 $1/q$ 成反比，
弹性约为 $-1$。矿井 1 的第二水源启动时刻由 3 min 延迟到 5 min 时，最晚充满时刻仅增加 1 min；
矿井 2 在 4--6 min 区间内最大指标不变。三档逃生速度倍率为 $0.9,1.0,1.1$ 时，矿井 1 每档三人中的最晚到达
依次为 $20.004,18.104,16.549\,\mathrm{min}$，矿井 2 依次为 $21.223,19.201,17.546\,\mathrm{min}$；
18 次扰动均到达原出口。安全阈值取 $0.25,0.30,0.35\,\mathrm{m}$ 时，六人的路线、出口和到达时刻均不变，
说明当前最优路线在该局部区间具有裕量，但不能外推到任意阈值。

\begin{figure}[H]
\centering
\includegraphics[width=0.92\textwidth]{sensitivity_summary.png}
\caption{流量、启动时刻、逃生速度和安全水深的单因素灵敏度}
\end{figure}

\section{模型评价}
\subsection{优点}
模型以累计体积保证质量守恒，允许延迟水源、汇流、相向回蓄和途中水态切换；
输入、事件、路径和 Excel 输出均可追溯。相较继续沿原路线，Q4 重规划最多节省 7.62 min，具有明确决策价值。

\subsection{不足与推广}
模型仍属于缺参数条件下的可解释近似：分叉采用平均分流，未包含糙率、局部阻力、压力波和真实水面坡度；
水平边的水源势定向抑制循环，但不能替代完整的非恒定明渠方程。若获得流量监测，可进一步用水力半径和坡降标定分流权重，
并将事件模型与有限体积法对照。当前 Q4 已从首次通知时刻起连续复演，不再依赖中途位置线性插值。
